\subsection{向量积分的性质}

命题 1. --- 映射

\[
(\mathbf {f}, \mu) \mapsto \int \mathbf {f} d \mu
\]

（\(\widetilde{\mathcal{K}} (\mathrm{X};\mathrm{E})\times \mathcal{M}(\mathrm{X};\mathbf{C})\) 到 \(\mathrm{E}^{\prime *}\) 中的）是双线性的。

该命题由 No.~1 的定义 1 立即得出。

设 \(\mathbf{u}\) 是 \(\mathrm{E}\) 到一个局部凸空间 \(\mathrm{F}\) 中的一个连续线性映射；我们知道转置 \(^t\mathbf{u}\) 是 \(\mathrm{F}'\) （\(\mathrm{F}\) 的对偶）到 \(\mathrm{E}'\) （\(\mathrm{E}\) 的对偶）中的一个线性映射；我们以 \(^{tt}\mathbf{u}\) 记 \(\mathrm{E}^{\prime *} \to \mathrm{F}^{\prime *}\) 的映射，即 \(^t\mathbf{u}\) 的（代数意义的）转置；当 \(\mathrm{E}\) 与 \(\mathrm{F}\) 都是 Hausdorff 的且分别典范地等同于 \(\mathrm{E}^{\prime *}\) 与 \(\mathrm{F}^{\prime *}\) 的子空间时， \(^{tt}\mathbf{u}\) 是映射 \(\mathbf{u}\) 的延拓。在这些记号下：

命题 2. --- 设 u 是 E 到一个局部凸空间 F 中的一个连续线性映射；对每个函数 \(\mathbf{f} \in \widetilde{\mathcal{K}}(\mathrm{X};\mathrm{E})\) ，函数 \(\mathbf{u} \circ \mathbf{f}\) 属于 \(\widetilde{\mathcal{K}}(\mathrm{X};\mathrm{F})\) ，且

\[
\int \mathbf {u} (\mathbf {f} (x)) d \mu (x) = ^ {t t} \mathbf {u} \left(\int \mathbf {f} (x) d \mu (x)\right).\tag{2}
\]

对每个连续线性形式 \(z^{\prime} \in F^{\prime}\) ，有 \(z^{\prime} \circ u \circ f = y^{\prime} \circ f\) （令 \(y^{\prime} = z^{\prime} \circ u = {}^{t}u(z^{\prime}) \in E^{\prime}\) ），由此得第一个断言；此外，对所有 \(z^{\prime} \in F^{\prime}\) ，

\[
\begin{array}{r l} \left\langle \int (\mathbf {u} \circ \mathbf {f}) d \mu , \mathbf {z} ^ {\prime} \right\rangle & = \int \langle \mathbf {u} \circ \mathbf {f}, \mathbf {z} ^ {\prime} \rangle d \mu = \int \langle \mathbf {f}, ^ {t} \mathbf {u} (\mathbf {z} ^ {\prime}) \rangle d \mu \\ & = \left\langle \int \mathbf {f} d \mu , ^ {t} \mathbf {u} (\mathbf {z} ^ {\prime}) \right\rangle = \left\langle {} ^ {t t} \mathbf {u} \left(\int \mathbf {f} d \mu\right), \mathbf {z} ^ {\prime} \right\rangle , \end{array}
\]

由此得公式 (2)。

命题 3. --- 对每个函数 \(g \in \mathcal{C}(\mathrm{X};\mathbf{C})\) 及每个函数 \(\mathbf{f} \in \widetilde{\mathcal{K}}(\mathrm{X};\mathrm{E})\) ，函数 \(g\mathbf{f}\) 属于 \(\widetilde{\mathcal{K}}(\mathrm{X};\mathrm{E})\) ，且

\[
\int \mathbf {f} d (g \cdot \mu) = \int \mathbf {f} g d \mu .\tag{3}
\]

事实上，对每个 \(z^{\prime} \in E^{\prime}\) ， \(\langle gf, z^{\prime} \rangle = g\langle f, z^{\prime} \rangle\) ，由此得第一个断言；此外，

\[
\begin{array}{r l} \left\langle \int \mathbf {f} d (g \cdot \mu), \mathbf {z} ^ {\prime} \right\rangle & = \int \langle \mathbf {f}, \mathbf {z} ^ {\prime} \rangle d (g \cdot \mu) = \int g \langle \mathbf {f}, \mathbf {z} ^ {\prime} \rangle d \mu \\ & = \int \langle g \mathbf {f}, \mathbf {z} ^ {\prime} \rangle d \mu = \left\langle \int g \mathbf {f} d \mu , \mathbf {z} ^ {\prime} \right\rangle , \end{array}
\]

由此得 (3)。

命题 4. --- 设 \(\mu\) 是 X 上的一个正测度，S 是其支集，f 是 \(\widetilde{\mathcal{K}}(\mathrm{X};\mathrm{E})\) 中的一个函数。设 E 是 Hausdorff 的，并给空间 \(\mathrm{E}^{\prime*}\) 赋以弱拓扑 \(\sigma(\mathrm{E}^{\prime*},\mathrm{E}^{\prime})\) 。

\begin{enumerate}
\def\labelenumi{(\roman{enumi})}
\item
  积分 \(\int \mathbf{f}d\mu\) 属于 C ，即 \(\mathrm{E}^{\prime *}\) 中由 \(\mathbf{f}(\mathrm{S})\) 生成的凸锥的闭包。
\item
  若 \(\mu\) 有界，则积分 \(\int \mathbf{f} d\mu\) 属于集 \(\| \mu \| \cdot D\) ，其中 \(D\) 是 \(\mathbf{f}(S)\) 在 \(\mathrm{E}'^*\) 中的闭凸包。
\end{enumerate}

若 \(\mathrm{E}\) 是复的，我们给 \(\mathrm{E}\) 赋以其底层实向量空间结构，如前所述，这并不改变公式 (1)。

\begin{enumerate}
\def\labelenumi{(\roman{enumi})}
\tightlist
\item
  我们知道 C 是 \(E^{\prime*}\) 中包含 \(\mathbf{f}(\mathrm{S})\) 且由过 0 的闭超平面确定的闭半空间之交；因此只需证明，对 \(z^{\prime} \in E^{\prime}\) ，关系 \(\langle\mathbf{f}(x), \mathbf{z}^{\prime}\rangle \geqslant 0\) 对一切 \(x \in S\) 成立蕴涵
\end{enumerate}

\[
\left\langle \int \mathbf {f} d \mu , \mathbf {z} ^ {\prime} \right\rangle \geqslant 0;
\]

但由于

\[
\left\langle \int \mathbf {f} d \mu , \mathbf {z} ^ {\prime} \right\rangle = \int \left\langle \mathbf {f}, \mathbf {z} ^ {\prime} \right\rangle d \mu ,
\]

这由 §2, No.~3, 命题 8 的推论 2 得出。

\begin{enumerate}
\def\labelenumi{(\roman{enumi})}
\setcounter{enumi}{1}
\tightlist
\item
  我们知道 D 是 \(E^{\prime*}\) 中包含 \(\mathbf{f}(S)\) 的闭半空间之交；因此只需证明，对 \(z^{\prime} \in E^{\prime}\) ，关系 \(\langle\mathbf{f}(x), \mathbf{z}^{\prime}\rangle \leqslant a\) 对一切 \(x \in S\) 成立蕴涵
\end{enumerate}

\[
\left\langle \int \mathbf {f} d \mu , \mathbf {z} ^ {\prime} \right\rangle \leqslant a \| \mu \|;
\]

而这由 §2, No.~3, 命题 8 的推论 3 得出。

推论. --- 设 \(\mu\) 是 X 上的一个正测度，f 是属于 \(\mathcal{K}(\mathrm{X};\mathrm{E})\) 的一个映射，D 是 \(\mathbf{f}(\mathrm{X})\) 在 \(\mathrm{E}'^*\) 中的闭凸包。则存在一个数 \(a > 0\) 使 \(\int \mathbf{f} d\mu \in a \cdot \mathrm{D}\) 。

若 \(\nu\) 是 \(X\) 上的任一测度，则存在数 \(a_1, a_2, a_3, a_4 > 0\) 使 \(\int \mathbf{f} d\nu \in a_1\mathrm{D} - a_2\mathrm{D} + ia_3\mathrm{D} - ia_4\mathrm{D}\) 。

先设 \(\mu\) 是正的；按假设，f 的支集 K 是紧的；若 \(\nu\) 是 \(\mu\) 在 K 的一个相对紧开邻域上的限制，则 \(\nu\) 有界，且由命题 4, (ii)， \(\int f d\mu = \int f d\nu \in \|\nu\| \cdot D\) 。第二个结果由此得出，因为任何复测度都可写成 \(\mu_{1} - \mu_{2} + i\mu_{3} - i\mu_{4}\) ，其中诸 \(\mu_{j}\) 是正测度。

命题 5. --- 设空间 X 是紧的，f 是 X 到一个 Hausdorff 局部凸空间 E 中的一个连续映射。 \(\mathbf{f}(\mathbf{X})\) 在 \(E^{\prime*}\) 中（对 \(\sigma(E^{\prime*}, E^{\prime})\) ）的闭凸包等于向量 \(\int f d\mu\) （对 X 上总质量为 1 的一切正测度 \(\mu\) 所取）所成之集。

设 C 是 \(\mathbf{f}(\mathbf{X})\) 在 \(E^{\prime*}\) 中的闭凸包；由于 \(\mathbf{f}(\mathbf{X})\) 是紧的且 \(E^{\prime*}\) 是完备的，故 C 是紧的。我们已经知道（命题 4），

\(\int \mathbf{f}d\mu \in \mathbb{C}\) 对属于 X 上总质量等于 1 的正测度所成凸集 H 的每个测度 \(\mu\) 成立。另一方面，H 对淡拓扑是凸的且紧的（§1, No.~9, 命题 15 的推论 3），且是质量为 1、支集有限的正测度所成凸集 \(\mathrm{H}_0\) 的（对该拓扑的）闭包（§2, No.~4, Th. 1 的推论 3）。现在， \(\mathrm{H}_0\) 在映射 \(\mu \mapsto \int \mathbf{f}d\mu\) 下的像是凸包 \(\mathrm{C}_0\) ，即 \(\mathbf{f}(\mathbf{X})\) 在 \(\mathrm{E}'^*\) 中的凸包。另一方面，这一映射对 \(\mathcal{M}(\mathrm{X};\mathbf{C})\) 上的淡拓扑与拓扑 \(\sigma (\mathrm{E}^{\prime *},\mathrm{E}^{\prime})\) （\(\mathrm{E}'^*\) 上的）是连续的，因为按定义 \(\langle \int \mathbf{f}d\mu ,\mathbf{z}'\rangle = \int \langle \mathbf{f},\mathbf{z}'\rangle d\mu\) ；于是 \(\mathrm{H} = \overline{\mathrm{H}}_0\) 的像是一个包含 \(\mathrm{C}_0\) 且含于 C 的紧凸集；由于 \(\mathrm{C} = \overline{\mathrm{C}}_0\) ，这个像等于 C。

命题 6. --- 对 X 到一个 Hausdorff 局部凸空间 E 中的每个具紧支集的连续映射 f、E 上每个连续半范数 q 及 X 上的每个测度 \(\mu\) （使 \(\int f d\mu \in E\) ），

\[
q \left(\int \mathbf {f} d \mu\right) \leqslant \int (q \circ \mathbf {f}) d | \mu |.\tag{4}
\]

设 D 是 \(z \in E\) （\(q(\mathbf{z}) \leqslant 1\) ）所成之集；D 是闭的、凸的且包含 0，因此 \(D = D^{\circ\circ}\) （TVS, II, §6, No.~3, Th. 1 的推论 3）。因此只需证明，对每个 \(z' \in D^{\circ}\) ，

\[
\left| \left\langle \int \mathbf {f} d \mu , \mathbf {z} ^ {\prime} \right\rangle \right| \leqslant \int (q \circ \mathbf {f}) d | \mu |;
\]

由于

\[
\left\langle \int \mathbf {f} d \mu , \mathbf {z} ^ {\prime} \right\rangle = \int \left\langle \mathbf {f}, \mathbf {z} ^ {\prime} \right\rangle d \mu ,
\]

且由于按 \(D^{\circ}\) 的定义，

\[
| \langle \mathbf {f} (x), \mathbf {z} ^ {\prime} \rangle | \leqslant q (\mathbf {f} (x))
\]

对所有 \(x \in \mathbf{X}\) 成立，故所欲证的不等式由 §1, No.~6 的不等式 (13) 得出。

